\begin{thebibliography}{99}
\bibitem{S} OpenAI, \emph{No bigeodesics in planar first-passage percolation}, stipulated source 212, repository revision \texttt{fd4aeeb2ee4fc729c18d98444fed42fd0529eeeb}. Exact source paths and blob identities are recorded in the accompanying source receipts.
\bibitem{AH} D. Ahlberg and C. Hoffman, \emph{Random coalescing geodesics in first-passage percolation}, \href{https://arxiv.org/abs/1609.02447v2}{arXiv:1609.02447v2}.
\bibitem{DH} M. Damron and J. Hanson, \emph{Busemann functions and infinite geodesics in two-dimensional first-passage percolation}, \href{https://arxiv.org/abs/1209.3036v1}{arXiv:1209.3036v1}; Communications in Mathematical Physics 325 (2014), 917--963, doi:10.1007/s00220-013-1875-y.
\bibitem{DHb} M. Damron and J. Hanson, \emph{Bigeodesics in first-passage percolation}, \href{https://arxiv.org/abs/1512.00804v2}{arXiv:1512.00804v2}, Conjecture 1.
\bibitem{ADH} A. Auffinger, M. Damron and J. Hanson, \emph{Limiting geodesics for first-passage percolation on subsets of $\Z^2$}, \href{https://arxiv.org/abs/1302.5413v2}{arXiv:1302.5413v2}; Annals of Applied Probability 25 (2015), 373--405, doi:10.1214/13-AAP999. Appendix C, especially Lemma C.5 and (26).
\bibitem{CT} R. Cerf and M. Th\'eret, \emph{Weak shape theorem in first passage percolation with infinite passage times}, \href{https://arxiv.org/abs/1404.4539v2}{arXiv:1404.4539v2}.
\end{thebibliography}
